\section{Hitting time mixing for large cycles}\label{sec: large cycles}

In this section, we prove \Cref{thm: large cycle}. Unless otherwise noted, we use the notations from \Cref{sec: Jain-Sawhney}. As we will see, the proof of \Cref{thm: large cycle} reduces to a Fourier estimate for the two-step distribution $P_n^{*2}$. The following proposition gives the required bound on the contribution from all nontrivial irreducible representations.

\begin{prop}\label{prop: fourier bound large cycle}
Suppose $n=\omega(1),1\le n-k\le o(n^{1/2})$.  Then, we have 
$$\sum_{\lambda\ne(n),(1^n)}\frac{\chi_\l(\tau_k)^4}{d_\lambda^2}=o(1).$$
\end{prop}

\begin{proof}[Proof of \Cref{thm: large cycle}, assuming \Cref{prop: fourier bound large cycle}]
After the first random cycle, the probability that we stop at the next cycle is exactly the probability that the second $k$-cycle covers the remaining $n-k$ cards. Therefore, we have
$$\mathbf{P}(\tau=2)=\frac{\binom{k}{n-k}}{\binom{n}{n-k}}=\left(1-\frac{n-k}{n}\right)\cdots\left(1-\frac{n-k}{k+1}\right)=1-O\left(\frac{(n-k)^2}{n}\right)=1-o(1).$$
Hence, it suffices to prove that $d_{TV}(P_n^{*2},U_{\A_n})=o(1)$. Recall from \eqref{eq: Fourier transform of Pn} that 
$$\widehat P_n^{*2}(\lambda)=(\widehat P_n(\lambda))^2=\frac{\chi_\lambda(\tau_k)^2}{d_\lambda^2}.$$
Moreover, the Fourier transform of $U_{\A_n}$ is one over the partitions $(n),(1^n)$, and zero on the other partitions. Now, applying the Cauchy-Schwarz inequality followed by the Plancherel formula, we have that
\begin{align}
\begin{split}
4\cdot d_{TV}(P_n^{*2},U_{\A_n})^2
&\le\sum_{\l\vdash n}d_\l\Tr\left((\widehat P_n^{*2}(\lambda)-\widehat U_{\A_n}(\l))(\widehat P_n^{*2}(\lambda)-\widehat U_{\A_n}(\l))^\dag\right)\\
&=\sum_{\lambda\ne(n),(1^n)}\frac{\chi_\l(\tau_k)^4}{d_\lambda^2}\\
&=o(1).\\
\end{split}
\end{align}
Therefore, we conclude $d_{TV}(P_n^{*2},U_{\A_n})=o(1)$, which completes the proof.
\end{proof}

\begin{rmk}
The assumption $k=n-o(n^{1/2})$ is sharp for the estimate in \Cref{prop: fourier bound large cycle}. Indeed, consider the partition $\lambda=(n-1,1)$. In this case,
$$
\chi_{(n-1,1)}(\tau_k)=d_{(n-k-1,1)}=n-k-1,
\qquad d_{(n-1,1)}=n-1.
$$
Therefore the single term corresponding to $\lambda=(n-1,1)$ in the spectral sum is
$$
\frac{\chi_{(n-1,1)}(\tau_k)^4}{d_{(n-1,1)}^2}=\frac{(n-k-1)^4}{(n-1)^2}.
$$
If $k=n-\Omega(n^{1/2})$, then this term is
already $\Omega(1)$. Thus the conclusion of the proposition cannot hold in that range, showing that the condition $k=n-o(n^{1/2})$ is optimal.
\end{rmk}

For the rest of this section, we aim to prove \Cref{prop: fourier bound large cycle}. We first record a consequence of the (recursive) Murnaghan-Nakayama rule for a single
large cycle. The uniqueness statement in the following proposition is the same observation as Lulov-Pak \cite[Lemma 5.1]{lulov2002rapidly}. For our later estimates, we also keep track of the resulting shape of this unique rim hook.

\begin{lemma}\label{lem: MN one term three types}
Let $n\ge 1$, let $k$ satisfy $1\le n-k<n/2$, and write $r:=n-k$.
Then for every partition $\lambda\vdash n$, 
there exists at most one $\tilde\lambda\vdash r$ such that $\lambda/\widetilde\lambda$ is a rim hook (also called border strip) of size $k$, and in this case we have 
$|\chi_\lambda(\tau_k)|=d_{\widetilde\lambda}$. Otherwise, we have $\chi_\lambda(\tau_k)=0$.
\end{lemma}

\begin{proof}
By the Murnaghan-Nakayama rule applied to the cycle type $(k,1^{n-k})$, we have
$$
\chi_\lambda(\tau_k)=\sum_{\substack{\widetilde\lambda\vdash n-k\\
\lambda/\widetilde\lambda\text{ is a rim hook of size }k}}
(-1)^{\operatorname{ht}(\lambda/\widetilde\lambda)-1}
d_{\widetilde\lambda}.
$$
Since $k>n/2$, Lulov-Pak \cite[Lemma 5.1]{lulov2002rapidly} shows that there is at most one way to remove a rim hook of size $k$ from $\lambda$. Hence the above sum has at most one term, and therefore 
$$|\chi_\lambda(\tau_k)|=d_{\widetilde\lambda}$$
when such a rim hook exists, while $\chi_\lambda(\tau_k)=0$ otherwise.
\end{proof}

\begin{prop}
\label{prop: fixed core rim hook sum}
Suppose that $r=o(n^{1/2})$, and fix a partition
$\widetilde\lambda\vdash r$ satisfying
$\widetilde\lambda_1'\le \widetilde\lambda_1\le r-1$. Then
$$
\sum_{\substack{\lambda\vdash n:\ \widetilde\lambda\subseteq\lambda\\
\lambda/\widetilde\lambda\text{ is a rim hook of size }n-r}}
\left(\frac{d_{\widetilde\lambda}}{d_\lambda}\right)^2
=
O\left(
\left(\frac{3r}{n}\right)^{2(r-\widetilde\lambda_1)}
\right).
$$
\end{prop}

\begin{proof}
Since $\lambda/\widetilde\lambda$ is a rim hook of size $n-r$, it is
contained in the outer rim of $\lambda$, and hence
$$
n-r\le \lambda_1+\lambda_1'-1.
$$
Since $r=o(n^{1/2})$, every partition appearing in the sum satisfies, for all sufficiently large $n$, either
$\lambda_1\ge \frac n3$ or $\lambda_1'\ge \frac n3$. We first consider the contribution of the partitions satisfying $\lambda_1\ge n/3$.

We will repeatedly use the following consequence of the hook-length
formula. Given a partition $\eta\vdash N$, we will write $\eta^*$ for the partition obtained from $\eta$ by deleting its first row, and we will Denote $s:=|\eta^*|$. Then
$$
\frac{d_\eta}{d_{\eta^*}}
=
\binom{N}{s}
\prod_{j=1}^{\eta_2}
\frac{\eta_1-j+1}
{\eta_1-j+1+\eta_j'-1}.
$$
If $\eta_1\ge n/3$ and $s\le r+1$, then $r=o(n^{1/2})$ implies
$$
\prod_{j=1}^{\eta_2}
\frac{\eta_1-j+1}
{\eta_1-j+1+\eta_j'-1}
=1-o(1),
$$
uniformly over all such $\eta$. Consequently,
$d_\eta\ge(1-o(1))\binom{N}{s}d_{\eta^*}$.
We will also use the upper bound from
\Cref{lem: dimension bound}, which claims that
$$
d_{\widetilde\lambda}
\le
\binom{r}{r-\widetilde\lambda_1}
d_{\widetilde\lambda^*}.
$$

To begin with, consider the partition $\lambda$ for which $\lambda/\widetilde\lambda$ is a horizontal strip added entirely to the first row. In this case,
$$
\lambda^*=\widetilde\lambda^*,
\qquad
|\lambda^*|=r-\widetilde\lambda_1.
$$
Applying the hook-length comparison to $\lambda$, together with
\Cref{lem: dimension bound} for $\widetilde\lambda$, gives
$$\frac{d_{\widetilde\lambda}}{d_\lambda}\le\frac{
\binom{r}{r-\widetilde\lambda_1}
d_{\widetilde\lambda^*}
}{(1-o(1))\binom{n}{r-\widetilde\lambda_1}
d_{\lambda^*}}=O\left(
\left(\frac{r}{n}\right)^{r-\widetilde\lambda_1}
\right).$$

For every other partition $\lambda$ with $\lambda_1\ge n/3$, trim the
rim hook immediately after its first downward turn, and denote the
remaining partition by $\mu$. Then $\mu\subseteq\lambda$, and the
diagram $\mu^*$ obtained by deleting the first row of $\mu$ contains
$\widetilde\lambda^*$ together with at least one additional box. In
particular,
$$
r-\widetilde\lambda_1+1\le |\mu^*|\le r.
$$
By the branching rule, $d_{\mu^*}\ge d_{\widetilde\lambda^*}$. Therefore, the hook-length formula and
\Cref{lem: dimension bound} give
$$\frac{d_{\widetilde\lambda}}{d_\lambda}\le\frac{d_{\widetilde\lambda}}{d_\mu}\le\frac{\binom{r}{r-\widetilde\lambda_1}d_{\widetilde\lambda^*}
}{(1-o(1))\binom{|\mu|}{|\mu^*|}d_{\mu^*}}\le(1+o(1))\frac{\binom{r}{r-\widetilde\lambda_1}}{\binom{|\mu|}{r-\widetilde\lambda_1+1}}=O\left(\left(\frac{3r}{n}\right)^{r-\widetilde\lambda_1+1}\right).$$
Here, we use the assumption that $|\mu|\ge \lambda_1\ge n/3$.

There are only $O(n)$ possible partitions $\lambda$ of this second
kind. Indeed, with $\widetilde\lambda$ fixed, one may imagine sliding
the rim hook along the outer boundary, starting with its horizontal
part in the top row and moving its turning point successively
downward until its vertical part reaches the bottom of the first
column. The position of this turning point determines $\lambda$, and
there are at most $O(n)$ possible positions.

It follows that
\begin{align*}
\sum_{\substack{\lambda\supseteq\widetilde\lambda\\
\lambda/\widetilde\lambda\text{ is a rim hook of size }n-r\\
\lambda_1\ge n/3}}
\left(\frac{d_{\widetilde\lambda}}{d_\lambda}\right)^2 &\le O\left(\left(\frac{r}{n}\right)^{
2(r-\widetilde\lambda_1)}\right)+O(n)
O\left(\left(\frac{3r}{n}\right)^{
2(r-\widetilde\lambda_1+1)}\right) \\
&=O\left(\left(\frac{3r}{n}\right)^{
2(r-\widetilde\lambda_1)}\right),
\end{align*}
where the last step follows from $r=o(n^{1/2})$.

Finally, consider the partitions satisfying $\lambda_1'\ge n/3$.
Applying the preceding argument after conjugation gives the bound
$$
O\left(
\left(\frac{3r}{n}\right)^{
2(r-\widetilde\lambda_1')}
\right).
$$
Since $\widetilde\lambda_1\ge\widetilde\lambda_1'$, we have $r-\widetilde\lambda_1'\ge r-\widetilde\lambda_1$. As $3r/n<1$ for all sufficiently large $n$, this contribution is also
$$
O\left(
\left(\frac{3r}{n}\right)^{
2(r-\widetilde\lambda_1)}
\right).
$$
Combining the two cases completes the proof.
\end{proof}

\begin{prop}
\label{prop: one row and one column cores}
Suppose that $r=o(n^{1/2})$. Then
$$
\sum_{\substack{\lambda\vdash n,\ \lambda\neq(n),(1^n)\\
\lambda/(r)\text{ is a rim hook of size }n-r}}
\frac{1}{d_\lambda^2}=O\left(\frac{r^2}{n}\right)=o(1).
$$
Similarly,
$$
\sum_{\substack{\lambda\vdash n,\ \lambda\neq(n),(1^n)\\
\lambda/(1^r)\text{ is a rim hook of size }n-r}}\frac{1}{d_\lambda^2}=O\left(\frac{r^2}{n}\right)=o(1).
$$
\end{prop}

\begin{proof}
The first estimate follows routinely from the proof of \Cref{prop: fixed core rim hook sum}. Indeed, when $\widetilde\lambda=(r)$, the exceptional horizontal extension is precisely $\lambda=(n)$, which is excluded from the sum. Every remaining partition belongs to the second case considered there. There are at most $O(n)$ such partitions, and each of them satisfies
$$
\frac{d_{(r)}}{d_\lambda}=\frac{1}{d_\lambda}
=O\left(\frac{2r}{n}\right).
$$
Consequently,
$$
\sum_{\substack{\lambda\vdash n,\ \lambda\neq(n),(1^n)\\
\lambda/(r)\text{ is a rim hook of size }n-r}}
\frac{1}{d_\lambda^2}
\le O(n)O\left(\frac{r^2}{n^2}\right)=O\left(\frac{r^2}{n}\right)=o(1).
$$

The second estimate follows from the first one by conjugation, using
$d_{\lambda'}=d_\lambda$.
\end{proof}

Finally, we are prepared for our proof of \Cref{prop: fourier bound large cycle}.

\begin{proof}[Proof of \Cref{prop: fourier bound large cycle}]
Write $r:=n-k$. By \Cref{lem: MN one term three types}, for every
$\lambda\vdash n$, either $\chi_\lambda(\tau_k)=0$, or there exists a
unique partition $\widetilde\lambda\vdash r$ such that
$\lambda/\widetilde\lambda$ is a rim hook of size $k$, in which case
$$
|\chi_\lambda(\tau_k)|=d_{\widetilde\lambda}.
$$
Therefore,
\begin{align*}
\sum_{\lambda\neq(n),(1^n)}
\frac{\chi_\lambda(\tau_k)^4}{d_\lambda^2}
&=
\sum_{\widetilde\lambda\vdash r}
d_{\widetilde\lambda}^2
\sum_{\substack{\lambda\vdash n:\,
\lambda/\widetilde\lambda\text{ is a rim hook of size }k\\
\lambda\neq(n),(1^n)}}
\left(\frac{d_{\widetilde\lambda}}{d_\lambda}\right)^2.
\end{align*}

Since conjugation preserves dimensions and sends the pair
$(\lambda,\widetilde\lambda)$ to
$(\lambda',\widetilde\lambda')$, it is enough, up to a factor of $2$,
to sum over the partitions satisfying
$\widetilde\lambda_1\ge\widetilde\lambda_1'$.
The contribution of $\widetilde\lambda=(r)$ is
$o(1)$ by \Cref{prop: one row and one column cores}. It remains to consider
$\widetilde\lambda\neq(r)$. For
$1\le l\le r-1$, \Cref{lem: dimension bound} gives
$$
d_{\widetilde\lambda}
\le
\binom{r}{r-l}d_{\widetilde\lambda^*}
$$
whenever $\widetilde\lambda_1=l$, where
$\widetilde\lambda^*$ is obtained by deleting the first row. Consequently, using the identity
$$
\sum_{\mu\vdash r-l}d_\mu^2=(r-l)!,
$$
we obtain
$$d_{\widetilde\lambda}^2\le
\binom{r}{r-l}^2
\sum_{\mu\vdash r-l}d_\mu^2=
\binom{r}{r-l}^2(r-l)!\le
\frac{r^{2(r-l)}}{(r-l)!}.$$

Applying \Cref{prop: fixed core rim hook sum} and then summing over
$l$ therefore gives
\begin{align*}
\sum_{\substack{\widetilde\lambda\vdash r,\,
\widetilde\lambda\neq(r)\\
\widetilde\lambda_1\ge\widetilde\lambda_1'}}
d_{\widetilde\lambda}^2
\sum_{\substack{\lambda\vdash n:\,
\lambda/\widetilde\lambda\text{ is a rim hook of size }k}}
\left(\frac{d_{\widetilde\lambda}}{d_\lambda}\right)^2 &\le
\sum_{l=1}^{r-1}
O\left(\left(\frac{3r}{n}\right)^{2(r-l)}
\right)\sum_{\substack{\widetilde\lambda\vdash r\\\widetilde\lambda_1=l}}
d_{\widetilde\lambda}^2 \\
&\le
\sum_{l=1}^{r-1}
O\left(
\frac{(3r^2/n)^{2(r-l)}}{(r-l)!}
\right) \\
&\le
O\left(
\exp\left(\left(\frac{3r^2}{n}\right)^2\right)-1
\right)\\
&=o(1),
\end{align*}
where the last equality follows from $r=o(n^{1/2})$.

The conjugate partitions contribute the same amount. Combining the
preceding estimates proves
$$
\sum_{\lambda\neq(n),(1^n)}
\frac{\chi_\lambda(\tau_k)^4}{d_\lambda^2}
=o(1).
$$
\end{proof}

